\section{Background and Related Work}

\label{sec:background}
% ══════════════════════════════════════════════════════════════════════════════

\subsection{Classical Species Distribution Models}

The dominant paradigm in SDM is the correlative approach, which relates
species presence (and sometimes absence) records to environmental
predictors. MaxEnt \citep{phillips2006maximum} maximizes entropy of the
predicted distribution subject to constraints from observed feature
means. Boosted regression trees \citep{elith2008working} and random
forests \citep{cutler2007random} offer non-parametric alternatives that
accommodate interactions among predictors. Ensemble approaches such as
BIOMOD2 \citep{thuiller2009biomod} combine multiple algorithms to reduce
model uncertainty.

Despite their utility, correlative SDMs rest on the assumption that
current species-environment relationships are stable under future
conditions---the \emph{stationarity assumption}---which climate change
directly violates \citep{dormann2012correlation}.

\subsection{Deep Learning for SDMs}

Neural networks were first applied to SDMs by \citet{lek1996application},
but modern deep learning architectures have dramatically expanded
capability. Convolutional neural networks (CNNs) applied to
remotely-sensed imagery can predict species occurrence from raw satellite
data \citep{deneu2021convolutional}. Deep species distribution models
(DeepSDM) use multi-layer perceptrons with environmental covariates as
inputs \citep{botella2018deep}. However, these approaches treat each
spatial location independently and do not leverage landscape structure.

\subsection{Graph Neural Networks}

GNNs generalize deep learning to graph-structured data
\citep{kipf2017semi, velickovic2018graph}. Message-passing neural
networks \citep{gilmer2017neural} iteratively update node representations
by aggregating information from neighbors. Spatial applications of GNNs
include traffic forecasting \citep{li2018diffusion}, molecular property
prediction \citep{duvenaud2015convolutional}, and social network analysis
\citep{hamilton2017inductive}. In ecology, GNNs remain underexplored,
though recent work has applied them to food web modeling
\citep{strydom2021graph} and metapopulation dynamics
\citep{mcinerny2022graph}.

\subsection{Connectivity and Corridor Modeling}

Circuit theory \citep{mcrae2006isolation, mcrae2008using} models
landscape connectivity as electrical resistance, where current flow
between nodes indicates movement probability. Least-cost path analysis
\citep{adriaensen2003application} identifies single optimal routes.
Omniscape \citep{landau2021omnidirectional} extends circuit theory to
continuous surfaces. These approaches require pre-specified resistance
surfaces, which are typically hand-crafted from land cover maps---a
limitation HabitatNet addresses by learning resistance from data.


% ══════════════════════════════════════════════════════════════════════════════
